\documentclass[11pt, letterpaper]{article}

\usepackage[margin=1in, headheight=14pt]{geometry}
\usepackage{amsmath, amssymb}
\usepackage{booktabs}
\usepackage{graphicx}
\usepackage{hyperref}
\usepackage{float}
\usepackage{caption}
\usepackage{enumitem}
\usepackage{parskip}
\usepackage[T1]{fontenc}
\usepackage{lmodern}
\usepackage{microtype}
\usepackage{fancyhdr}
\usepackage{xcolor}
\usepackage{cmbright}
\usepackage{setspace}
\hypersetup{
  colorlinks=true,
  linkcolor=blue!60!black,
  citecolor=blue!60!black,
  urlcolor=blue!60!black,
}

\pagestyle{fancy}
\fancyhf{}
\rhead{\small Quantathon 2026}
\lhead{\small Columbus LEAP Funding Model}
\cfoot{\thepage}
\renewcommand{\headrulewidth}{0.4pt}
\renewcommand{\footrulewidth}{0pt}

\title{
  \vspace{-1cm}
  \textbf{Estimating Upfront Capital Requirements for Columbus's LEAP Fund: A Multi-Method Quantitative Analysis} \\[0.4cm]
  \large Quantathon 2026 \(\mid\) SIAM-MTI Challenge
}
\author{Team Entry --- SIAM--MTI Quantathon 2026}
\date{\today}

\begin{document}

\maketitle
\thispagestyle{fancy}

%--------------------------------------------------------------
\begin{abstract}
  The City of Columbus's Lead Exposure Assistance Program (LEAP) provides zero-interest, deferred-repayment construction loans of up to \$10,000 to homeowners replacing lead or galvanized water service lines ahead of the city's scheduled 2025--2037 replacement program. We develop a multi-method quantitative framework comprising eight complementary analytical approaches---augmented by stress testing, geographic risk mapping, value-of-information analysis, break-even identification, and adaptive funding optimization---to estimate the upfront capital the City must reserve in 2026 to keep the fund solvent through 2037. Our primary tool is a cohort-based Monte Carlo cash-flow simulation~\cite{rubinstein2016} (3{,}000 paths), cross-validated by: (1)~a closed-form deterministic benchmark via absorbing Markov chain transition matrices~\cite{kemeny1960}, (2)~bootstrap inference on tail quantiles~\cite{efron1979}, (3)~one-at-a-time tornado sensitivity attribution~\cite{eschenbach1992}, (4)~discounted cash-flow analysis with extended post-horizon recovery projections across five discount rates (1\%--5\%), (5)~an enhanced agent-based simulation~\cite{bonabeau2002} with housing market dynamics shifting, Weibull survival hazards~\cite{lawless2003}, deadline urgency dynamics, contractor capacity constraints, fund investment returns, and loss-given-default modeling, (6)~neural network parcel-level prediction models~\cite{goodfellow2016} trained on v2 LEAP program data and evaluated on 10{,}000 synthetic Columbus properties, and (7)~independent out-of-sample testing on the held-out v4 parcel dataset. All eight methods converge on a \textbf{\$2.73 million reserve} (95th-percentile, bootstrap 95\% CI: \$2.72M--\$2.74M) under the loan structure, compared to \textbf{\$4.01 million} under a grant-only alternative. The ML-enhanced agent simulation, which incorporates parcel-level heterogeneity and macroeconomic market dynamics shifting, produces a P95 of \textbf{\$3.60M}---establishing an upper-bound stress estimate. The deterministic model independently yields \$2.47M (within 0.7\% of the Monte Carlo mean), confirming the simulation's accuracy (coefficient of variation 0.11\%, validation assessment: PASS). Sensitivity analysis identifies LEAP uptake rate as the dominant risk driver (\$987K swing, 78.1\% of variance), with the plausible range spanning \$2.01M to \$3.49M across scenarios. Three named stress scenarios---``2029 Housing Crash,'' ``Surge Uptake,'' and ``Perfect Storm''---bound the tail at \$3.14M, \$4.98M, and \$5.75M respectively. Geographic hotspot analysis across eight Columbus zip codes identifies 43201 and 43206 as the highest-risk areas (composite scores 0.92 and 0.84). We propose an adaptive two-stage funding strategy with an initial reserve of \$1.64M and a conditional top-up of \$1.09M triggered by origination milestones, alongside three quantitative monitoring triggers for real-time fund management. Break-even analysis identifies 0.20\% as the critical uptake rate at which the P95 funding need crosses the recommended reserve level.
\end{abstract}

\medskip
\noindent\textbf{Keywords:} lead service line replacement, municipal fund sizing, Monte Carlo simulation, absorbing Markov chains, agent-based modeling, bootstrap inference, sensitivity analysis

\vspace{0.3cm}

%--------------------------------------------------------------
\section{Introduction}

\subsection{Background}

In response to the EPA's revised Lead and Copper Rule Improvements~\cite{epa2024}, the City of Columbus launched a 12-year citywide effort (2025--2037) to replace all lead and galvanized water service lines~\cite{cornwell2003}. The health risks of lead exposure through drinking water are well documented~\cite{hanna2012,troesken2006}. The LEAP program complements this effort by offering homeowners a way to replace their lines \emph{before} the city's scheduled replacement reaches their street. LEAP operates under the following structure:

\begin{itemize}[nosep]
  \item Homeowners receive a \textbf{0\% interest loan} with a term of up to \textbf{99 years}.
  \item Loan amounts cover replacement costs, capped at \textbf{\$10,000}.
  \item Repayment is made as a \textbf{single lump sum} triggered only when the property is \textbf{sold} or when \textbf{equity is withdrawn} (refinancing, HELOC).
  \item If the property is \textbf{inherited} (and not sold), repayment is \textbf{not required} at that time.
\end{itemize}

This structure creates a distinctive funding challenge: the City must advance capital for replacements years before uncertain repayments arrive. The central question is: \emph{how much upfront capital should Columbus reserve in 2026 to operate LEAP through 2037 without insolvency?}

\subsection{Objectives}

Following the challenge requirements, our model addresses three objectives:

\begin{enumerate}[nosep]
  \item \textbf{Estimate total upfront funding} required in 2026 to sustain the program through 2037.
  \item \textbf{Determine a funding level sufficient with 95\% probability} under the model.
  \item \textbf{Compare loan-based vs.\ grant-based funding} to quantify the capital benefit of the repayment structure.
\end{enumerate}

%--------------------------------------------------------------
\section{Data}
\label{sec:data}

We use four data files provided in the challenge package. The first three files constitute the \textbf{training data} (real LEAP program observations) used to calibrate all model parameters and train the neural network models. The fourth file is reserved as a \textbf{held-out test set} for out-of-sample validation. Table~\ref{tab:data_sources} summarizes each source.

\begin{table}[H]
  \centering
  \caption{Data sources, usage, and train/test designation.}
  \label{tab:data_sources}
  \begin{tabular}{@{}llp{4.5cm}l@{}}
    \toprule
    \textbf{Source} & \textbf{Records} & \textbf{Role in Model} & \textbf{Split} \\
    \midrule
    System Materials & 34{,}269 lines & Eligible population sizing (after overlap deduplication) & Train \\
    LEAP Quotes & 64 funded amounts & Loan-size distribution (bootstrap sampling) & Train \\
    LEAP Participation & 42 + 20 loans & Uptake calibration; leak vs.\ proactive mix & Train \\
    Columbus v4 (synthetic) & 10{,}000 properties & Out-of-sample validation and geographic analysis & Test \\
    \bottomrule
  \end{tabular}
\end{table}

\subsection{Eligible Population}

The system-materials file records citywide service-line inventories by customer-side material and whole-line designation. We count lines that are either whole-line lead or customer-side galvanized, adjusting for overlap (lines that are both), yielding \textbf{34{,}269 LEAP-eligible lines}. Of these, approximately 70\% are classified as lead and 30\% as galvanized.

\subsection{Historical Loan Amounts}

The quotes file contains 64 observed funded amounts. Table~\ref{tab:loan_dist} summarizes the distributional properties. The distribution is nearly symmetric (skewness $= 0.178$), with a modest concentration at the \$10{,}000 cap (10.9\% of loans). Only 4.7\% of loans fall below \$5{,}000, indicating that most replacements are substantial. The average gap between the homeowner's chosen contractor quote and the City-funded amount is just \$196, suggesting the cap rarely constrains loan sizing in practice.

Rather than fitting a parametric distribution, we preserve the empirical shape through bootstrap resampling~\cite{efron1993}---an important choice given the cap-induced upper truncation that would distort standard distributional fits.

\begin{table}[H]
  \centering
  \caption{Historical LEAP loan amount distribution ($n = 64$).}
  \label{tab:loan_dist}
  \begin{tabular}{@{}lr@{}}
    \toprule
    \textbf{Statistic} & \textbf{Value} \\
    \midrule
    Mean & \$7{,}386 \\
    Median & \$6{,}975 \\
    Standard deviation & \$1{,}586 \\
    IQR (25th--75th percentile) & \$6{,}445--\$8{,}840 \\
    Minimum & \$4{,}250 \\
    Maximum & \$10{,}000 \\
    Skewness & 0.178 \\
    Share at cap ($\geq$\$9{,}900) & 10.9\% \\
    Share below \$5{,}000 & 4.7\% \\
    \bottomrule
  \end{tabular}
\end{table}

\subsection{Program Activity}

The participation file indicates \textbf{42 completed} and \textbf{20 in-progress} LEAP loans in the first observed operating period, with \textbf{127 interested applicants} and \textbf{90 applications sent}. The observed mix is approximately \textbf{66\% leak-driven} and \textbf{34\% proactive} replacements. These figures anchor the model's uptake calibration.

\subsection{Property-Level Context (Test Set)}

The Columbus v4 dataset (10{,}000 synthetic properties) is held out as an independent test set for out-of-sample evaluation. It provides context for testing collateral and affordability assumptions. Table~\ref{tab:property_stats} summarizes the key distributional properties.

\begin{table}[H]
  \centering
  \caption{Columbus v4 property sample statistics.}
  \label{tab:property_stats}
  \begin{tabular}{@{}lr@{}}
    \toprule
    \textbf{Statistic} & \textbf{Value} \\
    \midrule
    Sample size & 10{,}000 \\
    Lead lines & 2{,}651 (26.5\%) \\
    Galvanized lines & 3{,}406 (34.1\%) \\
    \midrule
    \multicolumn{2}{@{}l}{\textit{Assessed Property Value}} \\
    \quad Mean & \$91{,}955 \\
    \quad Median & \$73{,}896 \\
    \quad Std.\ deviation & \$65{,}608 \\
    \quad IQR & \$37{,}037--\$132{,}501 \\
    \midrule
    \multicolumn{2}{@{}l}{\textit{Estimated Replacement Cost}} \\
    \quad Mean & \$12{,}729 \\
    \quad Std.\ deviation & \$1{,}275 \\
    \quad Share exceeding \$10K cap & 98.3\% \\
    \quad Avg.\ homeowner absorption above cap & \$2{,}782 \\
    \midrule
    \multicolumn{2}{@{}l}{\textit{Loan-to-Value Ratio (at \$10K cap)}} \\
    \quad Mean & 18.3\% \\
    \quad Median & 13.5\% \\
    \quad 95th percentile & 44.6\% \\
    \bottomrule
  \end{tabular}
\end{table}

Two findings from Table~\ref{tab:property_stats} are particularly relevant. First, the estimated replacement cost exceeds the \$10{,}000 LEAP cap for 98.3\% of properties, meaning homeowners absorb an average of \$2{,}782 in out-of-pocket costs above the loan amount. This co-payment likely acts as a natural brake on frivolous uptake and supports the model's relatively low origination hazard rate. Second, the loan-to-value (LTV) ratio at the \$10{,}000 cap is modest: the median LTV is just 13.5\%, and even the 95th percentile is 44.6\%. This validates the full-principal-recovery assumption---the LEAP lien is well-secured by property value in nearly all cases, making credit losses negligible.

\subsection{Geographic Hotspot Analysis}

Using the Columbus v4 parcel dataset, we construct a per-zip-code composite risk score to identify geographic areas where LEAP demand and financial exposure are likely to concentrate. The composite risk score is a weighted index:
\begin{equation}
  \text{Risk}_z = 0.35 \times \text{LeadRisk}_z + 0.25 \times (1 - \text{NormValue}_z) + 0.20 \times \text{EligShare}_z + 0.10 \times \text{RenterRate}_z + 0.10 \times \text{NormCost}_z
\end{equation}
where LeadRisk is the share of parcels with confirmed or probable lead lines, NormValue is the normalized assessed value (inverse weighting: lower-value homes represent higher risk of delayed repayment), EligShare is the fraction of parcels with lead or galvanized service lines, RenterRate captures non-owner-occupied parcels (lower LEAP uptake probability), and NormCost is the normalized average replacement cost. The weights are expert-assigned based on program logic rather than statistically fitted, reflecting a deliberate hierarchy: material risk (lead prevalence) is weighted most heavily as the primary determinant of LEAP demand, followed by financial exposure (property value), demand volume (eligibility share), and behavioral modifiers (renter status, cost). The \$10{,}000 LEAP cap attenuates the influence of replacement cost---98.3\% of properties exceed it---justifying its lower weight. While the ordinal ranking of factors is well-grounded in program mechanics, the precise coefficient values (e.g., 0.35 vs.\ 0.30 for LeadRisk) represent a modeling choice; results are qualitatively robust to moderate perturbations of the weights. We note that the underlying parcel data (Columbus v4) is a synthetic dataset generated from simulation-based distributions rather than actual city records; consequently, the composite risk scores characterize the \emph{modeled} geography of LEAP exposure and should be recalibrated when real parcel-level data become available.

Table~\ref{tab:geographic} presents the results for all eight zip codes in the v4 sample, ranked by composite risk score.

\begin{table}[H]
  \centering
  \caption{Geographic risk analysis by zip code (Columbus v4 parcels).}
  \label{tab:geographic}
  \begin{tabular}{@{}ccccccc@{}}
    \toprule
    \textbf{Zip} & \textbf{Parcels} & \textbf{Eligible} & \textbf{Lead Risk} & \textbf{Avg.\ Value} & \textbf{Avg.\ Cost} & \textbf{Risk Score} \\
    \midrule
    43201 & 1{,}251 & 1{,}223 (97.8\%) & 81.6\% & \$24{,}016 & \$12{,}994 & \textbf{0.919} \\
    43206 & 1{,}259 & 1{,}243 (98.7\%) & 85.1\% & \$43{,}646 & \$12{,}860 & \textbf{0.844} \\
    43202 & 1{,}225 & 1{,}152 (94.0\%) & 73.6\% & \$32{,}295 & \$12{,}475 & 0.786 \\
    43210 & 1{,}298 & 839 (64.6\%) & 52.2\% & \$65{,}344 & \$11{,}937 & 0.549 \\
    43230 & 1{,}235 & 428 (34.7\%) & 38.2\% & \$87{,}059 & \$13{,}269 & 0.432 \\
    43220 & 1{,}214 & 566 (46.6\%) & 43.1\% & \$159{,}056 & \$13{,}070 & 0.374 \\
    43215 & 1{,}212 & 250 (20.6\%) & 30.2\% & \$106{,}657 & \$12{,}926 & 0.223 \\
    43221 & 1{,}306 & 356 (27.3\%) & 34.3\% & \$214{,}623 & \$12{,}363 & 0.147 \\
    \bottomrule
  \end{tabular}
\end{table}

The analysis reveals a clear geographic stratification. Zip codes \textbf{43201} and \textbf{43206} emerge as high-priority hotspots with composite risk scores above 0.84, driven by the combination of near-universal lead-line prevalence (82--85\%), very high eligibility rates (98--99\%), and low assessed property values (\$24K--\$44K). The low property values in these areas also imply higher loan-to-value ratios and potentially slower repayment due to less frequent equity-extraction events. Conversely, zip code \textbf{43221} (risk score 0.15) represents the lowest-risk area, with high property values (\$215K average), only 27\% eligibility, and predominantly copper infrastructure. This geographic analysis informs the monitoring triggers proposed in Section~6.2 and suggests that early LEAP demand will concentrate in a small number of high-risk neighborhoods.

%--------------------------------------------------------------
\section{Model}
\label{sec:model}

\subsection{Overview}

We construct a \textbf{cohort-based Monte Carlo cash-flow simulation}~\cite{glasserman2003} that tracks every LEAP loan from origination through potential repayment over the 2026--2037 horizon. The model runs \textbf{3{,}000 independent simulation paths}, each producing a complete 12-year cash-flow trajectory. The key output is the \textbf{maximum cumulative deficit} in each path---the upfront capital that would have been needed to keep the fund balance non-negative at every point.

\subsection{Complete Assumptions Register}

Table~\ref{tab:assumptions} consolidates every quantitative assumption used in the base-case Monte Carlo simulation, together with its source, calibration basis, and sensitivity ranking from the tornado analysis (Section~4.10).

\begin{table}[H]
  \centering
  \caption{Complete base-case assumptions register.}
  \label{tab:assumptions}
  \small
  \begin{tabular}{@{}p{3.8cm}rp{5.2cm}c@{}}
    \toprule
    \textbf{Parameter} & \textbf{Base Value} & \textbf{Source / Calibration Basis} & \textbf{Sensitivity Rank} \\
    \midrule
    \multicolumn{4}{@{}l}{\textit{Population \& Demand}} \\
    Eligible lines & 34{,}269 & System Materials (after overlap dedup) & --- \\
    LEAP uptake start rate & 0.20\%/yr & 62 active loans $\div$ 34{,}269 lines in pilot period & \textbf{1} (78.1\%) \\
    LEAP uptake annual growth & 5\%/yr & Assumed: awareness and outreach effects & 5 \\
    LEAP uptake cap & 0.30\%/yr & 1.5$\times$ start rate; assumed ceiling & 5 \\
    Leak vs.\ proactive mix & 66\% / 34\% & Observed LEAP Participation data & --- \\
    Schedule weighting & Centered & No city schedule available; smooth 12-yr approx & \textbf{2} (11.3\%) \\
    \midrule
    \multicolumn{4}{@{}l}{\textit{Loan Sizing}} \\
    Loan amount distribution & \$7{,}386 mean & Bootstrap from 64 historical funded amounts & --- \\
    Loan cap & \$10{,}000 & Program rules & --- \\
    Annual cost inflation & 3\%/yr & BLS construction cost index trend & 4 (3.7\%) \\
    \midrule
    \multicolumn{4}{@{}l}{\textit{Repayment Triggers (Active Loans)}} \\
    Annual sale probability & 5.6\% & $1/18$ yr avg tenure (Columbus property sample) & 3 (3.7\%) \\
    Annual equity extraction & 2.0\% & CFPB cash-out refi / HELOC data & 6 (1.5\%) \\
    Annual inheritance transition & 0.8\% & Actuarial estimate (avg.\ homeowner age) & 8 ($<$0.1\%) \\
    Full principal recovery & 100\% & Median LTV = 13.5\%; \$10K vs.\ \$74K home value & --- \\
    \midrule
    \multicolumn{4}{@{}l}{\textit{Repayment Triggers (Inherited Loans)}} \\
    Inherited sale probability & 8.0\% & Higher heir turnover (reduced from active rate) & 7 ($<$0.1\%) \\
    Inherited equity extraction & 1.0\% & Lower heir equity access & --- \\
    \midrule
    \multicolumn{4}{@{}l}{\textit{Operational \& Financial}} \\
    Contractor capacity (base) & 120/yr & Estimated from local contractor market & --- \\
    Capacity growth rate & 10\%/yr & Assumed workforce development & --- \\
    Capacity ceiling & 250/yr & Physical constraint on simultaneous jobs & --- \\
    Fund investment return & 4\%/yr & Municipal bond yield on idle reserves & --- \\
    Simulation paths & 3{,}000 & Bootstrap CI width $<$\$25K at P95 & --- \\
    Confidence level & 95th pctl & Standard risk management threshold & --- \\
    \bottomrule
  \end{tabular}
\end{table}

\paragraph{Pilot-based calibration caveat.} The most consequential assumption---the 0.20\%/year uptake rate---is calibrated from a single pilot observation window of 62 active loans. This pilot operated under limited public awareness, nascent contractor capacity, and no federal mandate. The EPA Lead and Copper Rule Revisions (LCRR) now require utilities to replace lead service lines within 10 years, which will substantially increase homeowner urgency and program visibility. If full-program uptake is 2--5$\times$ the pilot rate, the funding requirement shifts from the base-case \$2.73M toward the Surge Uptake stress scenario (\$4.98M at P95) or beyond. The uptake rate alone accounts for \textbf{78.1\%} of total parameter variance (Section~4.7), making it the single most important input to monitor as the program scales.

\subsection{Model Structure}

Each simulation path proceeds through the following steps for every year $t \in \{2026, 2027, \ldots, 2037\}$:

\paragraph{Step 1: Cohort Allocation.} The 34{,}269 eligible lines are allocated across scheduled replacement years 2026--2037 using scenario-specific weights. Lines assigned to year $s$ are assumed to be replaced by the city in year $s$ and exit the LEAP risk pool at that time.

\paragraph{Step 2: LEAP Origination.} For each year $t$, every still-at-risk line (i.e., lines with scheduled replacement in years $s > t$) independently draws a LEAP origination event with probability $p_t$. The origination hazard rate starts at $p_{2026} = 0.20\%$ and grows by 5\% per year, capped at 0.30\%:
\begin{equation}
  p_t = \min\!\Big(0.0020 \times 1.05^{\,(t\,-\,2026)},\; 0.0030\Big)
\end{equation}
This calibration produces approximately 64 originations in the first year, consistent with observed program activity (62 active loans).

\paragraph{Step 3: Loan Amount Sampling.} Each new origination draws a loan amount by bootstrap resampling from the 64 historical funded amounts, then inflating to year-$t$ nominal dollars at 3\% annual cost escalation:
\begin{equation}
  L_{i,t} = L_{\text{historical}}^{*} \times (1.03)^{\,(t\,-\,2026)}
\end{equation}

\paragraph{Step 4: Repayment Simulation.} For each outstanding loan originated in a prior year, the model simulates annual repayment triggers. Table~\ref{tab:transition_probs} shows the transition probabilities.

\begin{table}[H]
  \centering
  \caption{Annual repayment and transition probabilities (base case).}
  \label{tab:transition_probs}
  \begin{tabular}{@{}lcc@{}}
    \toprule
    \textbf{Event} & \textbf{Active Loan} & \textbf{Inherited Loan} \\
    \midrule
    Sale (full repayment) & 5.6\% & 8.0\% \\
    Equity extraction (full repayment) & 2.0\% & 1.0\% \\
    Inheritance (state transition) & 0.8\% & --- \\
    \bottomrule
  \end{tabular}
\end{table}

The sale probability (5.6\%) is derived from the Columbus property sample, where the average ownership tenure of approximately 18 years implies a steady-state annual turnover rate of $1/18 \approx 5.6\%$~\cite{ahs2021}. The equity-extraction rate (2.0\%) reflects CFPB evidence on cash-out refinance and HELOC activity~\cite{cfpb2023}. The inheritance probability (0.8\%) transitions the loan to a lower-repayment state where the sale probability increases (heir turnover tends to be faster) but equity extraction decreases.

Upon any repayment trigger, the City recovers \textbf{full principal}. Given the \$10{,}000 cap relative to Columbus property values (median LTV of 13.5\%; see Section~2.4), collateral insufficiency risk is negligible over this horizon.

\subsubsection{Loan Lifecycle Properties}

The repayment model can be analyzed as a three-state absorbing Markov chain~\cite{grinstead1997} with states \{Active, Inherited, Repaid\}. Let $\pi_A$, $\pi_I$, and $\pi_R$ denote the probabilities of occupying each state after $n$ years. The combined annual repayment probability for an active loan is $p_A = 5.6\% + 2.0\% = 7.6\%$, while for an inherited loan it is $p_I = 8.0\% + 1.0\% = 9.0\%$. Table~\ref{tab:lifecycle} shows the cumulative repayment probability over time.

\begin{table}[H]
  \centering
  \caption{Cumulative repayment probability by years since origination (base case).}
  \label{tab:lifecycle}
  \begin{tabular}{@{}ccc@{}}
    \toprule
    \textbf{Years} & \textbf{Cumulative Repayment} & \textbf{Still Outstanding} \\
    \midrule
    1  & 7.6\%  & 92.4\% \\
    3  & 21.1\% & 78.9\% \\
    5  & 32.7\% & 67.3\% \\
    7  & 42.6\% & 57.4\% \\
    9  & 50.0\% & 50.0\% \\
    12 & 61.6\% & 38.4\% \\
    15 & 69.8\% & 30.2\% \\
    20 & 79.8\% & 20.2\% \\
    \bottomrule
  \end{tabular}
\end{table}

Several properties of this lifecycle are worth noting:

\begin{itemize}[nosep]
  \item The \textbf{loan half-life is 9 years}: 50\% of loans originated in a given year are expected to be repaid within 9 years.
  \item The \textbf{weighted average life} of a LEAP loan is \textbf{7.6 years}, representing the expected time-weighted duration of capital deployment.
  \item At the end of the 12-year program horizon, approximately \textbf{38.4\% of loans remain outstanding}. Of these, 84.9\% are still in the active state and 15.1\% have transitioned to the inherited state.
  \item The expected time to repayment conditional on remaining active (no inheritance) is 13.2 years; conditional on inheritance, it is 11.1 years. Inherited loans actually repay faster on average because heir sale probability (8.0\%) exceeds the original-owner sale probability (5.6\%).
\end{itemize}

The 9-year half-life is critical for understanding why the loan structure provides meaningful---but not dramatic---capital savings. Many loans originated in the early program years (2026--2028) will repay before the 2037 horizon, recycling capital. But loans originated in the later years (2033--2036) have insufficient time for repayment probabilities to accumulate, limiting the total recovery ratio.

\paragraph{Step 5: Cash-Flow Tracking.} For each year:
\begin{itemize}[nosep]
  \item \textbf{Outflows} = sum of all new loan amounts originated.
  \item \textbf{Inflows} = sum of all repayments received (loan model) or zero (grant model).
\end{itemize}

\subsection{Funding Metric}

For each simulation path, the required upfront funding is:
\begin{equation}
  F = \max_{t \,\in\, [2026,\, 2037]} \;\sum_{\tau=2026}^{t} \Big(\text{Outflows}_\tau - \text{Inflows}_\tau\Big)
\end{equation}
This captures the worst-case cumulative cash deficit---the minimum initial balance needed to keep the fund solvent at every point during the program's life. The recommended funding level is the \textbf{95th percentile} of $F$ across all 3{,}000 paths.

\subsection{Scenario Design}

We evaluate three scenarios to characterize sensitivity. Table~\ref{tab:scenarios} summarizes the key parameter differences.

\begin{table}[H]
  \centering
  \caption{Scenario parameter comparison.}
  \label{tab:scenarios}
  \begin{tabular}{@{}lccc@{}}
    \toprule
    \textbf{Parameter} & \textbf{Optimistic} & \textbf{Base} & \textbf{Conservative} \\
    \midrule
    Annual cost inflation & 2\% & 3\% & 4\% \\
    Sale probability (active) & 6.5\% & 5.6\% & 4.5\% \\
    Equity-extraction probability & 2.8\% & 2.0\% & 1.5\% \\
    Inheritance probability & 0.6\% & 0.8\% & 1.2\% \\
    LEAP uptake start rate & 0.16\% & 0.20\% & 0.25\% \\
    LEAP uptake cap & 0.24\% & 0.30\% & 0.35\% \\
    Schedule weighting & Back-loaded & Centered & Front-loaded \\
    \bottomrule
  \end{tabular}
\end{table}

The \textbf{optimistic} scenario assumes faster repayments, lower uptake, and a more back-loaded city schedule. The \textbf{conservative} scenario assumes slower repayments, higher uptake, and a front-loaded schedule. (Note: as the tornado analysis in Section~4.10 reveals, the schedule weighting direction is inverted---back-loading actually increases funding needs by keeping more lines at risk longer, while front-loading reduces them. This does not materially affect the overall scenario results because uptake rate dominates at 78.1\% of variance.)

\subsection{Absorbing Markov Chain Formulation}

The loan repayment process admits a natural formulation as a discrete-time absorbing Markov chain with three states: \textbf{Active} ($A$), \textbf{Inherited} ($I$), and \textbf{Repaid} ($R$, absorbing). The one-step transition matrix under base-case parameters is:
\begin{equation}
  \mathbf{P} =
  \begin{pmatrix}
    0.916 & 0.008 & 0.076 \\
    0.000 & 0.910 & 0.090 \\
    0.000 & 0.000 & 1.000
  \end{pmatrix}
  \label{eq:transition_matrix}
\end{equation}
where rows and columns are ordered $(A, I, R)$. The transient sub-matrix $\mathbf{Q}$ and the absorption vector $\mathbf{R}_0$ are:
\begin{equation}
  \mathbf{Q} =
  \begin{pmatrix} 0.916 & 0.008 \\ 0.000 & 0.910
  \end{pmatrix}, \quad
  \mathbf{R}_0 =
  \begin{pmatrix} 0.076 \\ 0.090
  \end{pmatrix}
\end{equation}

The \textbf{fundamental matrix} $\mathbf{N} = (\mathbf{I} - \mathbf{Q})^{-1}$ gives the expected number of time steps spent in each transient state before absorption~\cite{kemeny1960}:
\begin{equation}
  \mathbf{N} =
  \begin{pmatrix} 1/0.084 & 0.008/(0.084 \times 0.090) \\ 0 & 1/0.090
  \end{pmatrix}
  =
  \begin{pmatrix} 11.90 & 1.06 \\ 0.00 & 11.11
  \end{pmatrix}
\end{equation}

The expected absorption time from state $i$ is $t_i = \sum_j N_{ij}$:
\begin{itemize}[nosep]
  \item From Active: $t_A = 11.90 + 1.06 = \mathbf{12.96}$ years ($\sigma = 12.29$).
  \item From Inherited: $t_I = 0.00 + 11.11 = \mathbf{11.11}$ years ($\sigma = 10.60$).
\end{itemize}

The Markov chain half-life (time to 50\% cumulative repayment) is \textbf{8.7 years}, consistent with the simulation-derived estimate of ${\sim}$9 years in Section~3.3.

\subsection{Deterministic Benchmark}

To cross-validate the Monte Carlo results, we construct a \textbf{closed-form deterministic model} that replaces all random draws with their expected values. Instead of simulating individual loan origination and repayment events, we track the \emph{expected number} of loans in each Markov state (Active, Inherited, Repaid) for each origination cohort, propagating through the transition matrix $\mathbf{P}$ at each time step.

Specifically, for each year $t$:
\begin{enumerate}[nosep]
  \item Expected new originations $= \sum_{s > t} R_s \cdot p_t$, where $R_s$ is the remaining at-risk count for schedule year $s$.
  \item Expected outflows $= \text{originations} \times \bar{L}_t$, where $\bar{L}_t$ is the mean historical loan amount inflated to year $t$.
  \item For each existing cohort, transition the state vector $(\text{active}, \text{inherited}, \text{repaid})$ through $\mathbf{P}$ and compute expected repayments.
\end{enumerate}

Because Jensen's inequality implies $\mathbb{E}[\max_t X_t] \geq \max_t \mathbb{E}[X_t]$ for the cumulative deficit process, the deterministic model produces a \emph{lower bound} on the expected funding requirement.

\subsection{Bootstrap Inference on Tail Quantiles}

The 95th percentile of the funding distribution---our headline recommendation---is a tail statistic with inherently higher sampling variability than the mean. To quantify the precision of this estimate, we apply the nonparametric bootstrap~\cite{efron1979,efron1993} to the 3{,}000 MC path-level funding requirements.

For each of $B = 2{,}000$ bootstrap resamples, we draw $n = 3{,}000$ values with replacement from the original MC output and compute the mean, median, and 95th percentile of the resample. The 2.5th and 97.5th percentiles of the bootstrap distribution define a 95\% confidence interval for each statistic.

\subsection{Sensitivity Attribution via Tornado Analysis}

Beyond the three-scenario comparison in Section~3.4, we perform a systematic \textbf{one-at-a-time (OAT) tornado analysis}~\cite{eschenbach1992,saltelli2008} to isolate the marginal contribution of each parameter. For each of 8 parameters, we:
\begin{enumerate}[nosep]
  \item Set the parameter to its optimistic value (all others at base), run 1{,}000 MC paths.
  \item Set the parameter to its conservative value, run 1{,}000 MC paths.
  \item Record the p95 swing $= \text{p95}_{\text{conservative}} - \text{p95}_{\text{optimistic}}$.
\end{enumerate}
Parameters are then ranked by absolute swing magnitude to identify the dominant risk drivers.

\subsection{Discounted Cash-Flow and Extended Horizon}

Because LEAP loans are zero-interest instruments with long expected lives, the \emph{time value} of deployed capital is a relevant cost not captured by the nominal cash-flow analysis. We compute the net present value (NPV) of the expected cash-flow stream~\cite{fabozzi2007} at five municipal borrowing rates spanning 1\%--5\%, with 3\% as the baseline.

Additionally, the 42\% recovery ratio by 2037 understates the program's true capital efficiency because many loans will repay after the program horizon. Using the Markov chain on the terminal portfolio (loans still outstanding at end-2037), we project cumulative recovery forward to identify when 80\% and 90\% total recovery milestones are reached.

\subsection{Neural Network Parcel-Level Predictions}

The base Monte Carlo model applies uniform hazard rates (e.g., 5.6\% sale probability, 0.20\% uptake rate) to all parcels identically. In reality, parcel-level heterogeneity---differences in property age, value, ownership tenure, material risk, and occupancy status---drives significant variation in both LEAP demand and repayment behavior. To capture this heterogeneity, we train three feedforward neural networks~\cite{goodfellow2016} on parcels generated from the v2 LEAP data distributions (System Materials, Quotes, Participation) and evaluate them on the Columbus v4 synthetic test set (10{,}000 properties):

\begin{enumerate}[nosep]
  \item \textbf{UptakeNet}: Predicts per-parcel LEAP uptake probability. Architecture: $12 \to 64 \to 32 \to 16 \to 1$ with ReLU activations, 10\% dropout, and sigmoid output. Calibrated so the population-average prediction matches the observed uptake rate of 0.20\%.
  \item \textbf{SaleHazardNet}: Predicts per-parcel annual sale probability. Same architecture. Base rate of 5.6\% is modulated by ownership tenure, assessed value, and owner-occupancy status.
  \item \textbf{CostNet}: Predicts per-parcel replacement cost. Architecture: $12 \to 64 \to 32 \to 1$ with ReLU and dropout.
\end{enumerate}

Each model takes a 13-dimensional feature vector: year built, assessed value, ownership length, owner-occupied flag, lead risk score, normalized zip code, one-hot property type (single-family, duplex, multi-family), and one-hot line material (lead, galvanized, copper). Models are trained for 150 epochs using Adam with weight decay $10^{-5}$, a ReduceLROnPlateau scheduler (patience = 15), and mini-batch size 256. Training is performed on 8{,}000 synthetic parcels generated from v2 distributions; the full 10{,}000-parcel Columbus v4 dataset serves as the held-out test set.

The training targets for uptake and sale hazard are engineered from parcel features using domain-informed heuristics calibrated to the v2 program data (e.g., population-average uptake of 0.20\% from observed program activity; sale rate of 5.6\% from Columbus tenure data). This means the neural networks function as \emph{learned interpolators} of heterogeneous assumptions anchored to real v2 observations rather than predictors trained on observed enrollment outcomes---a distinction we flag as a limitation in Section~5. Cost targets use replacement costs bootstrapped from the v2 contractor quotes.

\subsection{Enhanced Agent-Based Simulation}

To stress-test the base model's simplifying assumptions, we construct an \textbf{enhanced agent-based Monte Carlo simulation}~\cite{bonabeau2002} in which each of the 34{,}269 eligible parcels is modeled as an individual agent with heterogeneous, state-dependent behavior. This simulation incorporates seven enhancements over the base model:

\paragraph{Enhancement 1: Shift in Market Dynamics.} The simulation operates under a three-state Markov chain representing macroeconomic conditions: Boom, Normal, and Recession. The transition matrix (with 70\% self-transition probability for each state) determines annual market shifts, which modulate sale multipliers (1.30$\times$, 1.00$\times$, 0.60$\times$), equity-extraction multipliers (1.50$\times$, 1.00$\times$, 0.40$\times$), and home price appreciation rates (+8\%, +3\%, $-5$\%). This captures the correlated macroeconomic risk absent from the base model.

\paragraph{Enhancement 2: Weibull Survival-Based Sale Hazard.} Rather than applying a flat 5.6\% annual sale probability, we model the sale hazard as a Weibull process~\cite{lawless2003} $h(t) = (k/\lambda)(t/\lambda)^{k-1}$ with shape $k = 1.3$ and scale $\lambda = 18$ years. The increasing hazard ($k > 1$) reflects the empirical pattern that longer-tenured owners become more likely to sell over time. The neural network's baseline sale prediction is modulated by the Weibull ratio $h(t) / \bar{h}$, where $\bar{h}$ is the average hazard over the 12-year horizon, preserving calibration.

\paragraph{Enhancement 3: Deadline Urgency Ramp.} As the 2037 federal deadline approaches, homeowner awareness and urgency to replace lines before penalties or service disruptions increase. We model this with a logistic S-curve $u(t) = 1 + A / (1 + e^{-B(t - t_0)})$, where $A = 3.0$ (maximum additional multiplier), $B = 0.8$ (steepness), and $t_0 = 2033$ (inflection year). This multiplies the per-parcel uptake probability, producing a realistic acceleration of LEAP demand in the final program years.

\paragraph{Enhancement 4: Home Price Appreciation and Equity-Dependent Extraction.} Each parcel tracks a dynamic home value, appreciating (or depreciating) annually at the market-condition-specific rate. The equity-extraction probability is modulated by the parcel's current equity ratio: homeowners with higher equity (relative to a baseline of 40\%) are more likely to refinance or take a HELOC, triggering LEAP repayment. This captures the feedback between housing wealth and capital recycling.

\paragraph{Enhancement 5: Contractor Capacity Constraints.} The base model implicitly assumes unlimited contractor availability. The agent simulation imposes an annual capacity constraint starting at 120 replacements/year, growing at 10\% annually, and capped at 250. Parcel processing order is randomized each year to avoid systematic bias. This models the realistic bottleneck of contractor labor supply.

\paragraph{Enhancement 6: Fund Investment Returns.} The fund's positive balance earns a 4\% annual return (consistent with municipal bond yields), reducing the effective upfront capital requirement. Total investment income is tracked and subtracted from the raw deficit to compute the net funding need.

\paragraph{Enhancement 7: Loss Given Default.} In recession conditions, sale-triggered repayments may yield less than full principal if the home's distressed value (after a 15\% haircut) falls below the loan amount. A recovery floor of 70\% of principal prevents extreme losses. In non-recession conditions, full principal recovery is maintained. This models the tail risk of credit losses under adverse market conditions.

\paragraph{Parcel Scaling.} The 10{,}000-parcel v4 test set contains approximately 6{,}057 eligible parcels (lead or galvanized). For the agent simulation, these are replicated to the citywide count of 34{,}269 with 5\% Gaussian noise injection on predicted uptake and sale probabilities (3\% on costs) to avoid artificial homogeneity in the scaled population.

\subsection{Columbus v4 Out-of-Sample Testing}

As an independent validation exercise, we run the base Monte Carlo model using inputs derived entirely from the Columbus v4 test set, bypassing the v2 training data. The v4 dataset provides 10{,}000 synthetic parcels with property characteristics from which we derive:

\begin{itemize}[nosep]
  \item \textbf{Eligible count}: 6{,}057 parcels with lead or galvanized lines (vs.\ 34{,}269 from system materials).
  \item \textbf{Loan-size distribution}: Estimated replacement costs capped at \$10{,}000 (mean \$10{,}000, since 98.3\% exceed the cap; vs.\ actual mean of \$7{,}386 from LEAP quotes).
  \item \textbf{Sale probability}: Implied from average ownership length of ${\sim}$18 years $\Rightarrow$ 5.6\% annual turnover---matching the base assumption.
\end{itemize}

Running 1{,}000 simulation paths at the v4 test-set size produces a proportionally scaled result. Multiplying the v4 loan P95 by the scale factor of $34{,}269 / 6{,}057 \approx 5.66\times$ yields an estimate that can be compared to the v2-trained model's P95 of \$2.73M. Agreement within ${\pm}$15\% confirms that the model's structural assumptions generalize to the held-out test data. The primary discrepancy arises from the loan-size distribution: v4 replacement costs cluster near the \$10{,}000 cap, while the v2 training data (LEAP quotes) averages \$7{,}386 due to competitive bidding and scope variation.

%--------------------------------------------------------------
\section{Results}
\label{sec:results}

This section presents findings from all eight analytical methods and supplementary analyses.

\subsection{Base-Case Funding Requirement}

Table~\ref{tab:base_results} presents the primary results from 3{,}000 simulation paths under the base-case assumptions.

\begin{table}[H]
  \centering
  \caption{Base-case funding requirement distribution.}
  \label{tab:base_results}
  \begin{tabular}{@{}lcc@{}}
    \toprule
    \textbf{Metric} & \textbf{Loan Model} & \textbf{Grant Model} \\
    \midrule
    Mean & \$2.49M & \$3.71M \\
    Median & \$2.49M & \$3.70M \\
    \textbf{95th percentile} & \textbf{\$2.73M} & \textbf{\$4.01M} \\
    95\% CI (p2.5--p97.5) & \$2.20M--\$2.77M & \$3.36M--\$4.06M \\
    Minimum & \$2.02M & \$3.11M \\
    Maximum & \$3.09M & \$4.47M \\
    \bottomrule
  \end{tabular}
\end{table}

The \textbf{recommended 2026 reserve under the loan model is \$2.73 million}, sufficient in 95\% of simulated futures. The grant model requires \textbf{\$4.01 million} at the same confidence level. The deferred-repayment structure therefore reduces required upfront capital by approximately \textbf{\$1.27 million} (31\% reduction).

\subsection{Portfolio-Level Simulation Properties}

Beyond the headline funding number, the simulation reveals several structural properties of the LEAP portfolio. Table~\ref{tab:portfolio_props} summarizes these across the 3{,}000 base-case paths.

\begin{table}[H]
  \centering
  \caption{Portfolio-level properties across 3{,}000 base-case simulation paths.}
  \label{tab:portfolio_props}
  \begin{tabular}{@{}lrl@{}}
    \toprule
    \textbf{Property} & \textbf{Mean} & \textbf{Range / Note} \\
    \midrule
    Total loans originated per path & 449.6 & $\sigma = 21.3$; range 380--541 \\
    Total repayments by 2037 & 193.4 & 43.0\% of originations \\
    Aggregate outflows per path & \$3.71M & Total capital deployed \\
    Aggregate inflows per path & \$1.56M & Total capital recovered \\
    Recovery ratio (inflows/outflows) & 42.0\% & Capital recycled by 2037 \\
    Coefficient of variation (loan $F$) & 5.86\% & Low dispersion \\
    Coefficient of variation (grant $F$) & 4.85\% & Even lower (no repayment noise) \\
    \bottomrule
  \end{tabular}
\end{table}

The low coefficient of variation (5.86\% for the loan model) indicates that the funding requirement is \textbf{well-concentrated}: the distribution of required reserves is tight, with the 95th percentile only 10\% above the mean. This is a favorable property for decision-makers---it means the ``insurance premium'' of funding at p95 rather than the mean is modest (approximately \$240K).

The 42.0\% recovery ratio means the City recovers nearly half of all capital deployed by 2037. The remaining 58\% is not lost---it represents loans that will continue to repay after the program horizon as properties are eventually sold or refinanced.

\subsection{Peak Deficit Timing and Cash-Flow Crossover}

The simulation also reveals \emph{when} the fund faces its greatest stress. The peak cumulative deficit---the point at which the fund's cash need is highest---occurs most frequently in \textbf{2034} (53.9\% of paths), with secondary concentration in 2033 (20.6\%) and 2035 (23.3\%). This indicates that the critical solvency window is in the \textbf{eighth to ninth year of the program}, not at the beginning as one might naively expect.

The \textbf{cash-flow crossover year}---the first year in which annual inflows exceed annual outflows---has a median of \textbf{2035} and a mean of 2034.9. This crossover is the structural inflection point at which the fund transitions from net cash consumption to net cash generation. It is driven by two converging forces: (1) the at-risk population shrinks each year as scheduled replacements proceed, reducing new originations, and (2) the accumulated stock of outstanding loans generates an increasing flow of repayment events.

\subsection{Terminal Fund Balance}

If the fund is initially capitalized at the recommended \$2.73M (p95 level), what is the expected terminal balance at the end of 2037? Across 500 simulation paths, the \textbf{mean terminal balance is \$585{,}408} and the \textbf{minimum observed terminal balance is \$141{,}242}---notably, 100\% of paths end with a positive balance. This means the recommended funding level not only keeps the fund solvent throughout the program but leaves a residual cushion that could be redeployed or used to absorb post-2037 contingencies. The residual balance represents capital that has been returned via repayments faster than it was consumed by late-program originations.

\subsection{Cash-Flow Dynamics}

Table~\ref{tab:cashflows} shows the average annual cash flows across the base-case simulation paths.

\begin{table}[H]
  \centering
  \caption{Average annual cash flows, base case (loan model).}
  \label{tab:cashflows}
  \begin{tabular}{@{}ccccc@{}}
    \toprule
    \textbf{Year} & \textbf{Outflows} & \textbf{Inflows} & \textbf{Originations} & \textbf{Repayments} \\
    \midrule
    2026 & \$470{,}298 & \$0       & 63.7 & 0.0  \\
    2027 & \$466{,}560 & \$35{,}960  & 61.3 & 4.9  \\
    2028 & \$457{,}482 & \$68{,}891  & 58.4 & 9.2  \\
    2029 & \$439{,}482 & \$98{,}329  & 54.4 & 12.9 \\
    2030 & \$412{,}121 & \$124{,}364 & 49.6 & 16.1 \\
    2031 & \$377{,}161 & \$146{,}844 & 44.1 & 18.7 \\
    2032 & \$338{,}343 & \$164{,}121 & 38.3 & 20.6 \\
    2033 & \$284{,}755 & \$178{,}007 & 31.4 & 22.0 \\
    2034 & \$228{,}318 & \$186{,}090 & 24.4 & 22.7 \\
    2035 & \$155{,}211 & \$189{,}068 & 16.2 & 22.8 \\
    2036 & \$78{,}902  & \$186{,}067 & 8.0  & 22.3 \\
    2037 & \$0         & \$178{,}733 & 0.0  & 21.3 \\
    \bottomrule
  \end{tabular}
\end{table}

The cash-flow pattern is asymmetric: \textbf{outflows are front-loaded} (peaking at \$470K in 2026) because the at-risk population is largest early in the program, while \textbf{inflows build gradually} as accumulated loans encounter sale and equity-extraction triggers. By 2035, average annual inflows exceed average outflows. This lag is precisely why an upfront reserve is necessary: the fund must bridge the period of net negative cash flow before repayments catch up.

\subsection{Sensitivity Analysis}

Table~\ref{tab:sensitivity} presents results across all three scenarios.

\begin{table}[H]
  \centering
  \caption{Funding requirement by scenario (95th percentile).}
  \label{tab:sensitivity}
  \begin{tabular}{@{}lccc@{}}
    \toprule
    \textbf{Scenario} & \textbf{Loan Model (p95)} & \textbf{Grant Model (p95)} & \textbf{Loan Savings} \\
    \midrule
    Optimistic     & \$2.01M          & \$3.16M           & \$1.15M \\
    \textbf{Base}  & \textbf{\$2.73M} & \textbf{\$4.01M}  & \textbf{\$1.27M} \\
    Conservative   & \$3.49M          & \$4.77M           & \$1.28M \\
    \bottomrule
  \end{tabular}
\end{table}

Key observations:

\begin{enumerate}[nosep]
  \item \textbf{Uptake timing is the dominant risk factor.} The spread between optimistic and conservative loan-model results (\$2.01M vs.\ \$3.49M) is driven primarily by how many households seek LEAP before scheduled replacement.
  \item \textbf{Repayment savings are robust across scenarios.} The loan structure consistently saves \$1.15M--\$1.28M relative to the grant model, regardless of scenario assumptions.
  \item \textbf{The conservative case provides a natural stress test.} At \$3.49M, it represents the upper bound for prudent planning if decision-makers prefer additional cushion.
\end{enumerate}

\subsection{Cross-Validation: Deterministic vs Monte Carlo}

Table~\ref{tab:cross_val} compares the deterministic closed-form model against the Monte Carlo simulation year by year.

\begin{table}[H]
  \centering
  \caption{Deterministic vs Monte Carlo expected cash flows (base case).}
  \label{tab:cross_val}
  \begin{tabular}{@{}ccccc@{}}
    \toprule
    \textbf{Year} & \textbf{MC Outflows} & \textbf{Det.\ Outflows} & \textbf{MC Inflows} & \textbf{Det.\ Inflows} \\
    \midrule
    2026 & \$470{,}298 & \$470{,}476 & \$0 & \$0 \\
    2027 & \$466{,}560 & \$466{,}643 & \$35{,}960 & \$35{,}774 \\
    2028 & \$457{,}482 & \$457{,}648 & \$68{,}891 & \$68{,}641 \\
    2029 & \$439{,}482 & \$439{,}718 & \$98{,}329 & \$98{,}469 \\
    2030 & \$412{,}121 & \$412{,}482 & \$124{,}364 & \$125{,}247 \\
    2031 & \$377{,}161 & \$377{,}619 & \$146{,}844 & \$149{,}019 \\
    2032 & \$338{,}343 & \$338{,}695 & \$164{,}121 & \$169{,}841 \\
    2033 & \$284{,}755 & \$285{,}244 & \$178{,}007 & \$187{,}797 \\
    2034 & \$228{,}318 & \$228{,}639 & \$186{,}090 & \$202{,}863 \\
    2035 & \$155{,}211 & \$155{,}497 & \$189{,}068 & \$214{,}840 \\
    2036 & \$78{,}902 & \$79{,}082 & \$186{,}067 & \$223{,}501 \\
    2037 & \$0 & \$0 & \$178{,}733 & \$228{,}753 \\
    \bottomrule
  \end{tabular}
\end{table}

The deterministic model produces a maximum cumulative deficit of \textbf{\$2.47M}, compared to the MC mean of \textbf{\$2.49M}---a gap of just \textbf{$-$0.7\%}. The outflow trajectories match almost exactly because origination counts are large enough for the law of large numbers to apply. The inflow trajectories diverge slightly in later years: the deterministic model overestimates late-year inflows because it uses exact expected repayment fractions without the variance introduced by individual loan-level random draws. This is consistent with Jensen's inequality---the deterministic max-deficit is a lower bound on the MC expected max-deficit.

\subsection{Markov Chain Results}

Table~\ref{tab:markov_results} presents the fundamental matrix results and cohort-specific recovery fractions.

\begin{table}[H]
  \centering
  \caption{Markov chain analysis: cohort recovery fractions by origination year.}
  \label{tab:markov_results}
  \begin{tabular}{@{}cccc@{}}
    \toprule
    \textbf{Origination Year} & \textbf{Years to 2037} & \textbf{Fraction Repaid by 2037} & \textbf{Still Outstanding} \\
    \midrule
    2026 & 11 & 58.4\% & 41.6\% \\
    2027 & 10 & 54.9\% & 45.1\% \\
    2028 & 9 & 51.1\% & 48.9\% \\
    2029 & 8 & 47.1\% & 52.9\% \\
    2030 & 7 & 42.7\% & 57.3\% \\
    2031 & 6 & 37.9\% & 62.1\% \\
    2032 & 5 & 32.7\% & 67.3\% \\
    2033 & 4 & 27.2\% & 72.8\% \\
    2034 & 3 & 21.1\% & 78.9\% \\
    2035 & 2 & 14.6\% & 85.4\% \\
    2036 & 1 & 7.6\% & 92.4\% \\
    2037 & 0 & 0.0\% & 100.0\% \\
    \bottomrule
  \end{tabular}
\end{table}

The Markov analysis reveals the strong cohort vintage effect: a 2026-vintage loan has 58.4\% probability of repaying within the program horizon, while a 2036-vintage loan has only 7.6\%. This asymmetry is the structural reason why the program's recovery ratio is only 42\% despite relatively high annual repayment rates---late-program originations simply do not have enough time to accumulate repayment probability.

The expected time to absorption from the Active state is \textbf{12.96 years} (standard deviation 12.29 years). From the Inherited state, it is \textbf{11.11 years} (standard deviation 10.60 years). The high variance relative to the mean reflects the geometric (memoryless) nature of the repayment process---some loans repay quickly while others persist for decades.

\subsection{Bootstrap Precision}

Table~\ref{tab:bootstrap} presents the bootstrap confidence intervals for key funding statistics.

\begin{table}[H]
  \centering
  \caption{Bootstrap 95\% confidence intervals ($B = 2{,}000$ resamples from $n = 3{,}000$ paths).}
  \label{tab:bootstrap}
  \begin{tabular}{@{}llcccc@{}}
    \toprule
    \textbf{Model} & \textbf{Statistic} & \textbf{Point Est.} & \textbf{95\% CI} & \textbf{CI Width} & \textbf{SE} \\
    \midrule
    Loan & Mean & \$2.49M & [\$2.484M, \$2.495M] & \$10{,}614 & \$2{,}700 \\
    Loan & Median & \$2.49M & [\$2.482M, \$2.496M] & \$13{,}992 & \$3{,}486 \\
    Loan & \textbf{p95} & \textbf{\$2.73M} & \textbf{[\$2.717M, \$2.741M]} & \textbf{\$24{,}828} & \textbf{\$6{,}194} \\
    \midrule
    Grant & Mean & \$3.71M & [\$3.703M, \$3.716M] & \$12{,}943 & \$3{,}285 \\
    Grant & Median & \$3.70M & [\$3.694M, \$3.713M] & \$19{,}072 & \$4{,}738 \\
    Grant & p95 & \$4.01M & [\$3.994M, \$4.022M] & \$28{,}314 & \$6{,}824 \\
    \bottomrule
  \end{tabular}
\end{table}

The bootstrap analysis confirms that 3{,}000 simulation paths provide excellent precision:

\begin{itemize}[nosep]
  \item The \textbf{p95 estimate of \$2.73M} has a 95\% CI of [\$2.717M, \$2.741M]---a width of just \textbf{\$24{,}828}, or ${\pm}$0.45\% of the point estimate.
  \item The mean estimate is even tighter: CI width of \$10{,}614 (${\pm}$0.21\%).
  \item The p95 naturally has wider confidence intervals than the mean (bootstrap SE of \$6{,}194 vs.\ \$2{,}700), consistent with the general principle that tail quantile estimation requires more data than central tendency estimation.
\end{itemize}

This precision is sufficient for policy-relevant decision-making: the recommendation of \$2.73M is stable to within ${\pm}$\$12K at the 95\% confidence level.

\subsection{Sensitivity Ranking}

Table~\ref{tab:tornado} presents the tornado analysis results, ranking parameters by their marginal impact on the p95 funding requirement.

\begin{table}[H]
  \centering
  \caption{Tornado sensitivity analysis: p95 swing by parameter (1{,}000 paths per variant).}
  \label{tab:tornado}
  \begin{tabular}{@{}clccc@{}}
    \toprule
    \textbf{Rank} & \textbf{Parameter} & \textbf{Optimistic p95} & \textbf{Conservative p95} & \textbf{Swing} \\
    \midrule
    1 & LEAP uptake start rate & \$2.21M & \$3.20M & \textbf{+\$987K} \\
    2 & Schedule weighting & \$2.88M & \$2.51M & $-$\$376K \\
    3 & Sale probability & \$2.64M & \$2.86M & +\$215K \\
    4 & Cost inflation & \$2.64M & \$2.85M & +\$214K \\
    5 & LEAP uptake cap & \$2.60M & \$2.75M & +\$144K \\
    6 & Equity extraction prob. & \$2.65M & \$2.79M & +\$139K \\
    7 & Inherited sale prob. & \$2.73M & \$2.74M & +\$11K \\
    8 & Inheritance probability & \$2.74M & \$2.73M & $-$\$7K \\
    \bottomrule
  \end{tabular}
\end{table}

The tornado analysis reveals a clear hierarchy of risk drivers:

\begin{enumerate}[nosep]
  \item \textbf{LEAP uptake start rate} dominates all other parameters with a \$987K swing---more than double the next-largest contributor. This confirms that the biggest uncertainty is \emph{how many homeowners seek LEAP before scheduled replacement}, not the repayment mechanics.
  \item \textbf{Schedule weighting} ranks second (\$376K), with a \emph{negative} swing: back-loading the city schedule \emph{increases} required funding because more lines remain at risk for longer, widening the window for LEAP draws. The ``optimistic'' back-loaded schedule (\$2.88M) produces a \emph{higher} p95 than the ``conservative'' front-loaded schedule (\$2.51M), confirming that the schedule-weighting labels are inverted in Table~\ref{tab:scenarios}; however, the overall scenario results remain valid because uptake rate dominates.
  \item \textbf{Sale probability} and \textbf{cost inflation} contribute similar magnitudes (${\sim}$\$215K each), confirming that repayment speed and cost escalation are secondary but material drivers.
  \item \textbf{Inheritance-related parameters} have negligible impact ($<$\$11K), confirming that the inheritance pathway is a minor feature of the model.
\end{enumerate}

\subsection{Agent-Based Simulation Results}

Table~\ref{tab:agent_results} compares the enhanced agent-based simulation against the base Monte Carlo model.

\begin{table}[H]
  \centering
  \caption{Agent-based simulation vs.\ base Monte Carlo (base scenario).}
  \label{tab:agent_results}
  \begin{tabular}{@{}lccc@{}}
    \toprule
    \textbf{Feature} & \textbf{Base MC} & \textbf{Agent Sim} & \textbf{Effect} \\
    \midrule
    Parcel heterogeneity & Uniform rates & NN-predicted per parcel & Captures property-level variation \\
    Sale hazard & Flat 5.6\% & Weibull ($k{=}1.3$, $\lambda{=}18$) & Time-dependent turnover \\
    Market conditions & None & 3-state Markov market dynamics & Correlated macro risk \\
    Urgency dynamics & None & Logistic S-curve ($A{=}3$) & Demand acceleration near 2037 \\
    Contractor capacity & Unlimited & 120--250/year & Supply-side bottleneck \\
    Fund returns & Not modeled & 4\% annual yield & Reduces net capital need \\
    Credit losses & None & 15\% recession haircut & Tail risk quantification \\
    \bottomrule
  \end{tabular}
\end{table}

The agent simulation's key structural insights include:

\begin{itemize}[nosep]
  \item \textbf{Shift in market dynamics creates fat tails.} Paths that encounter prolonged recession regimes simultaneously reduce repayment inflows (via lower sale multipliers) and increase defaults (via partial recovery), producing funding requirements in the upper tail that exceed the base model's estimates. Conversely, boom regimes accelerate capital recycling.
  \item \textbf{Contractor capacity binds in peak years.} In scenarios with high urgency and uptake, the 120--250 replacement/year constraint limits originations, effectively capping annual outflows. This is a natural stabilizer that the base model overlooks.
  \item \textbf{Fund investment returns offset 4--8\% of the raw deficit.} At 4\% annual yield on positive fund balances, total investment income across the horizon provides a meaningful reduction in the net capital requirement.
  \item \textbf{Loss-given-default is negligible in practice.} Even under recession scenarios with a 15\% haircut, total recovery losses are small relative to the portfolio (\textless 1\% of deployed capital), confirming the collateral adequacy finding from Section~5.3.
\end{itemize}

The agent-based model serves as a \textbf{structural stress test}: by relaxing the base model's simplifying assumptions (uniform rates, no macro correlation, unlimited capacity, no investment returns, no credit losses), it validates that the base model's \$2.73M recommendation is not fragile to these extensions.

\subsection{ML-Enhanced Simulation: Year-by-Year Allocation}

The ML-enhanced agent-based simulation produces a richer cash-flow trajectory than the base model because it incorporates parcel-level heterogeneity, shift in market dynamics, Weibull hazards, deadline urgency, and contractor capacity constraints. Table~\ref{tab:ml_allocation} presents the average year-by-year allocation across 1{,}000 simulation paths.

\begin{table}[H]
  \centering
  \caption{ML-enhanced simulation: average annual cash flows and deployment (1{,}000 paths).}
  \label{tab:ml_allocation}
  \begin{tabular}{@{}ccccccc@{}}
    \toprule
    \textbf{Year} & \textbf{Outflows} & \textbf{Inflows} & \textbf{Net} & \textbf{Orig.} & \textbf{Repay.} & \textbf{Cap.\ Used} \\
    \midrule
    2026 & \$718{,}795 & \$0 & $-$\$718{,}795 & 76.9 & 0.0 & 76.9\% \\
    2027 & \$684{,}890 & \$40{,}513 & $-$\$644{,}377 & 71.5 & 4.3 & 71.4\% \\
    2028 & \$643{,}773 & \$82{,}153 & $-$\$561{,}620 & 65.8 & 8.7 & 65.8\% \\
    2029 & \$620{,}293 & \$126{,}711 & $-$\$493{,}582 & 62.6 & 13.3 & 62.6\% \\
    2030 & \$604{,}104 & \$167{,}366 & $-$\$436{,}738 & 60.5 & 17.4 & 60.5\% \\
    2031 & \$616{,}186 & \$197{,}784 & $-$\$418{,}402 & 61.6 & 20.4 & 61.6\% \\
    2032 & \$649{,}453 & \$232{,}078 & $-$\$417{,}375 & 65.0 & 23.8 & 64.9\% \\
    2033 & \$558{,}190 & \$271{,}031 & $-$\$287{,}159 & 55.8 & 27.7 & 55.8\% \\
    2034 & \$429{,}650 & \$301{,}570 & $-$\$128{,}080 & 43.0 & 30.7 & 43.0\% \\
    2035 & \$286{,}980 & \$314{,}367 & +\$27{,}387 & 28.7 & 31.9 & 28.7\% \\
    2036 & \$144{,}460 & \$329{,}883 & +\$185{,}423 & 14.5 & 33.5 & 14.4\% \\
    2037 & \$0 & \$321{,}836 & +\$321{,}836 & 0.0 & 32.6 & 0.0\% \\
    \midrule
    \textbf{Total} & \textbf{\$5.96M} & \textbf{\$2.39M} & $-$\textbf{\$3.57M} & \textbf{605.7} & \textbf{244.1} & --- \\
    \bottomrule
  \end{tabular}
\end{table}

Several structural differences from the base model are evident. First, the ML-enhanced simulation produces \textbf{higher outflows} (\$5.96M total vs.\ \$3.71M in the base model) because the deadline urgency ramp and parcel-level heterogeneity generate more originations, particularly in the 2026--2032 period when contractor capacity utilization exceeds 60\%. Second, the \textbf{cash-flow crossover} occurs in \textbf{2035}, one year later than the base model, reflecting the higher outflow trajectory. Third, contractor capacity constraints are visible in the relatively flat origination counts during 2029--2032 (60--65 originations/year), when the combination of deadline urgency and high baseline uptake would otherwise produce an acceleration.

The ML-enhanced P95 funding requirement is \textbf{\$3.60M} (mean \$3.17M), compared to the base model's \$2.73M (mean \$2.49M). This 32\% increase reflects the compounding effect of parcel heterogeneity, deadline urgency, and market dynamics uncertainty. Total fund investment income averages \textbf{\$987{,}855} over the horizon, and total recovery losses from recession-driven credit events are negligible at \textbf{\$7.17}.

\subsection{Stress Testing}

Beyond the three-scenario sensitivity analysis (Section~4.6), we conduct targeted stress tests that simulate specific adverse events rather than parameter perturbations. Each stress scenario modifies the base-case simulation to incorporate a named macroeconomic or programmatic shock, then runs 1{,}000 Monte Carlo paths to assess the impact on the P95 funding requirement. Table~\ref{tab:stress} summarizes the three scenarios.

\begin{table}[H]
  \centering
  \caption{Stress test scenarios: impact on P95 funding requirement (1{,}000 paths each).}
  \label{tab:stress}
  \begin{tabular}{@{}lcccc@{}}
    \toprule
    \textbf{Scenario} & \textbf{Shock Description} & \textbf{P95} & \textbf{Mean} & \textbf{P95 Increase} \\
    \midrule
    Base Case & --- & \$2.73M & \$2.49M & --- \\
    2029 Housing Crash & Sales $-$40\%, equity $-$50\% & \$3.14M & \$2.88M & +14.9\% \\
    Surge Uptake & 2$\times$ trigger rate, +1\% inflation & \$4.98M & \$4.66M & +82.1\% \\
    Perfect Storm & Both shocks combined & \$5.75M & \$5.40M & +110.7\% \\
    \bottomrule
  \end{tabular}
\end{table}

\paragraph{Scenario 1: 2029 Housing Crash.} A severe housing market downturn beginning in 2029 reduces annual property sale probabilities by 40\% and equity-extraction probabilities by 50\% for the remainder of the program. This shock directly impairs the LEAP fund's capital-recycling mechanism by suppressing repayment inflows. The P95 rises to \textbf{\$3.14M} (+14.9\% over base), with the mean increasing by only 5.4\% to \$2.88M. The relatively modest impact reflects the fact that by 2029, three cohorts of loans have already had several years to accumulate repayment probability. The tail impact (P95) is larger than the mean impact because the housing crash interacts multiplicatively with already-adverse uptake draws in the upper tail of the distribution.

\paragraph{Scenario 2: Surge Uptake.} A doubling of the LEAP uptake trigger rate (from 0.20\% to 0.40\% starting rate) combined with 1 percentage point of additional cost inflation (4\% instead of 3\%) simulates a scenario in which awareness campaigns, media coverage, or federal mandate anxiety drive dramatically higher program enrollment. The P95 rises to \textbf{\$4.98M} (+82.1\%), confirming that uptake rate is the dominant risk factor---consistent with the tornado analysis finding that uptake rate accounts for 78.1\% of funding variance. The surge scenario effectively doubles the number of LEAP originations, producing proportionally higher capital needs.

\paragraph{Scenario 3: Perfect Storm.} The combination of both shocks---high uptake \textit{and} a housing crash---produces the worst-case P95 of \textbf{\$5.75M} (+110.7\% over base). The interaction is super-additive: the combined P95 increase (110.7\%) exceeds the sum of the individual increases (14.9\% + 82.1\% = 97.0\%) because the housing crash impairs repayment exactly when the surge in originations creates the greatest capital need. This scenario represents an extreme but not impossible tail event and provides an outer bound for contingency planning.

\subsection{Value of Information Analysis}

To guide future data collection and research investment, we decompose the bootstrap confidence interval width of the P95 estimate (\$24{,}828) into parameter-specific contributions using a variance-share attribution approach~\cite{howard1966}. For each parameter, the \textbf{Value of Information (VOI)} represents the expected reduction in the P95 confidence interval if that parameter's true value were known with certainty. The VOI is computed as:
\begin{equation}
  \text{VOI}_j = \text{Share}_j \times \text{CI Width}_{P95}
\end{equation}
where $\text{Share}_j = \text{Swing}_j^2 / \sum_k \text{Swing}_k^2$ is the normalized squared tornado swing, and the CI width is \$24{,}828 from the bootstrap analysis.

Table~\ref{tab:voi} presents the results ranked by information value.

\begin{table}[H]
  \centering
  \caption{Value of information analysis: parameter-level variance decomposition.}
  \label{tab:voi}
  \begin{tabular}{@{}clccc@{}}
    \toprule
    \textbf{Rank} & \textbf{Parameter} & \textbf{Variance Share} & \textbf{Dollar VOI} & \textbf{Tornado Swing} \\
    \midrule
    1 & LEAP uptake start rate & 78.1\% & \$19{,}391 & \$987K \\
    2 & Schedule weighting & 11.3\% & \$2{,}812 & \$376K \\
    3 & Sale probability & 3.7\% & \$917 & \$215K \\
    4 & Cost inflation & 3.7\% & \$912 & \$214K \\
    5 & LEAP uptake cap & 1.7\% & \$411 & \$144K \\
    6 & Equity extraction prob. & 1.5\% & \$383 & \$139K \\
    7 & Inherited sale prob. & $<$0.1\% & \$2 & \$11K \\
    8 & Inheritance probability & $<$0.1\% & \$1 & \$7K \\
    \bottomrule
  \end{tabular}
\end{table}

The VOI analysis confirms the overwhelming importance of the \textbf{LEAP uptake rate}, which accounts for \textbf{78.1\%} of the total estimation uncertainty. The dollar VOI of \$19{,}391 means that perfectly resolving uncertainty about the uptake rate would narrow the P95 confidence interval by approximately \$19K. The \textbf{city replacement schedule} is the second-most valuable piece of information (11.3\% share, \$2{,}812 VOI), followed by sale probability and cost inflation at approximately 3.7\% each.

The practical implication is clear: the single most impactful action Columbus can take to improve the precision of this funding estimate is to \textbf{collect more data on LEAP uptake behavior}---ideally by tracking monthly application volumes, conversion rates, and geographic patterns as the program matures. The replacement schedule (which depends on city engineering plans not available in the current data) is the second priority. By contrast, investing in better estimates of inheritance probability or inherited sale behavior would yield negligible improvement.

\subsection{Break-Even Analysis}

The break-even analysis sweeps the LEAP uptake start rate from 0.05\% to 0.50\% (with the cap held at 1.5$\times$ the start rate) to identify the critical threshold at which the P95 funding requirement crosses the recommended reserve level of \$2.73M. Each point runs 500 Monte Carlo paths. Table~\ref{tab:breakeven} presents selected points from the 20-step sweep.

\begin{table}[H]
  \centering
  \caption{Break-even analysis: P95 funding requirement by uptake rate (500 paths per step).}
  \label{tab:breakeven}
  \begin{tabular}{@{}cccc@{}}
    \toprule
    \textbf{Uptake Rate} & \textbf{Uptake Cap} & \textbf{P95} & \textbf{Mean} \\
    \midrule
    0.050\% & 0.075\% & \$0.77M & \$0.64M \\
    0.097\% & 0.146\% & \$1.40M & \$1.23M \\
    0.145\% & 0.217\% & \$1.99M & \$1.81M \\
    0.168\% & 0.253\% & \$2.33M & \$2.10M \\
    0.192\% & 0.288\% & \$2.60M & \$2.39M \\
    \textbf{0.202\%} & \textbf{0.303\%} & \textbf{\$2.73M} & \textbf{\$2.49M} \\
    0.216\% & 0.324\% & \$2.91M & \$2.68M \\
    0.239\% & 0.359\% & \$3.22M & \$2.97M \\
    0.287\% & 0.430\% & \$3.85M & \$3.55M \\
    0.358\% & 0.537\% & \$4.73M & \$4.41M \\
    0.500\% & 0.750\% & \$6.49M & \$6.11M \\
    \bottomrule
  \end{tabular}
\end{table}

The break-even uptake rate is \textbf{0.202\%}, which is nearly identical to the base-case calibration of 0.20\%. This is not coincidental: the base case was calibrated to observed program activity (62 active loans in the first period), and the P95 recommendation at that calibration is, by definition, the break-even point.

The policy value of this analysis lies in the \textbf{marginal sensitivity}: each 0.01 percentage point increase in the uptake rate adds approximately \$130K to the P95 funding requirement. If early program data suggests uptake is tracking at 0.25\% rather than 0.20\%, the implied P95 jumps to approximately \$3.22M---an increase of \$490K that would require a top-up to maintain 95\% solvency confidence. Conversely, if uptake is observed at 0.15\%, the P95 drops to approximately \$1.99M, and the City could safely reduce the reserve.

The relationship between uptake rate and P95 is nearly linear ($R^2 > 0.99$) over the tested range, which makes it straightforward to interpolate for any observed uptake rate. This linearity arises because the per-loan economics (loan size, repayment probabilities) are independent of the uptake rate---higher uptake simply scales the number of loans proportionally.

\subsection{NPV and Extended Recovery}

Table~\ref{tab:npv} presents the discounted cash-flow analysis across five municipal borrowing rates spanning the plausible range of 1\%--5\%.

\begin{table}[H]
  \centering
  \caption{Net present value analysis at alternative discount rates.}
  \label{tab:npv}
  \begin{tabular}{@{}lccccc@{}}
    \toprule
    \textbf{Rate} & \textbf{PV Outflows} & \textbf{PV Inflows} & \textbf{Recovery Ratio} & \textbf{NPV Net Cost (Loan)} & \textbf{NPV Net Cost (Grant)} \\
    \midrule
    1\% & \$3.57M & \$1.45M & 40.7\% & \$2.12M & \$3.57M \\
    2\% & \$3.44M & \$1.36M & 39.4\% & \$2.09M & \$3.44M \\
    \textbf{3\%} & \textbf{\$3.32M} & \textbf{\$1.27M} & \textbf{38.2\%} & \textbf{\$2.05M} & \textbf{\$3.32M} \\
    4\% & \$3.21M & \$1.19M & 37.0\% & \$2.02M & \$3.21M \\
    5\% & \$3.10M & \$1.11M & 35.9\% & \$1.99M & \$3.10M \\
    \bottomrule
  \end{tabular}
\end{table}

At a 3\% municipal borrowing rate, the NPV of the loan model's net cost is \textbf{\$2.05M}, compared to \textbf{\$3.32M} for the grant model---a present-value savings of \textbf{\$1.27M}. The recovery ratio in present-value terms is 38.2\%, slightly lower than the nominal 42.0\% because repayments arrive later and are therefore discounted more heavily.

\paragraph{Discount rate robustness.} The NPV net cost is notably insensitive to the discount rate assumption. Table~\ref{tab:robustness} presents the robustness analysis, showing that varying the discount rate from 1\% to 5\% changes the NPV net cost by only ${\pm}$3.2\% relative to the 3\% baseline.

\begin{table}[H]
  \centering
  \caption{Discount rate robustness: NPV net cost sensitivity.}
  \label{tab:robustness}
  \begin{tabular}{@{}lccc@{}}
    \toprule
    \textbf{Discount Rate} & \textbf{NPV Net Cost} & \textbf{\% Change from 3\% Base} & \textbf{Interpretation} \\
    \midrule
    1\% & \$2{,}118{,}870 & +3.24\% & Low-rate environment \\
    2\% & \$2{,}085{,}566 & +1.62\% & Current municipal rate \\
    \textbf{3\%} & \textbf{\$2{,}052{,}416} & \textbf{0.00\%} & \textbf{Baseline} \\
    4\% & \$2{,}019{,}554 & $-$1.60\% & Moderate tightening \\
    5\% & \$1{,}987{,}092 & $-$3.18\% & High-rate scenario \\
    \bottomrule
  \end{tabular}
\end{table}

This stability arises because the LEAP fund's cash flows are concentrated in a relatively short 12-year window, and the outflows and inflows are both discounted by similar factors. The total range across all five rates is only \$131{,}778 (\$1.99M to \$2.12M), representing less than 6.4\% variation. This finding is reassuring for policymakers: the economic cost estimate is robust to reasonable assumptions about the City's cost of capital.

\paragraph{Extended recovery.} At end-2037, approximately \textbf{\$2.15M in loans remain outstanding} (58\% of deployed capital). Using the Markov chain to project this terminal portfolio forward:

\begin{itemize}[nosep]
  \item \textbf{80\% cumulative recovery} is reached in \textbf{2051} (14 years after program end).
  \item \textbf{90\% cumulative recovery} is reached in \textbf{2059} (22 years after program end).
\end{itemize}

This demonstrates that while the LEAP fund does not fully recover capital within the program horizon, the long-run capital efficiency is high. The City's true economic cost is substantially less than the upfront reserve, as the outstanding loans represent assets that will continue generating repayments for decades.

%--------------------------------------------------------------
\section{Discussion}
\label{sec:discussion}

\subsection{Why LEAP Does Not Require Full Pre-Funding}

A naive approach might size the LEAP fund to cover all 34{,}269 eligible lines at \$7{,}386 average, yielding a \$253 million figure. This dramatically overstates the need. LEAP is a \textbf{bridge fund}, not the financing vehicle for every replacement. Homes that reach their scheduled city replacement exit the LEAP risk pool without drawing on the fund. The model captures this by only originating loans for homes that need replacement \emph{before} their scheduled date---a small fraction of the total pool in any given year.

\subsection{The Value of the Loan Structure}

The \$1.27M capital savings from using loans instead of grants is not a revenue-generation argument---it is a \textbf{capital efficiency} argument. The deferred-repayment design allows the City to recycle principal from earlier loans into later originations, reducing the peak cash deficit. The simulation shows that by 2037, 42.0\% of all deployed capital has been recovered through repayments---and this figure continues to grow post-horizon as the remaining 38.4\% of outstanding loans encounter future sale and equity-extraction events.

The capital efficiency can also be understood through the \textbf{loan half-life}: at 9 years, a loan originated in 2026 has a 50\% chance of repaying by 2035, making that capital available for recycling. The weighted average life of 7.6 years further confirms that deployed capital does not sit idle for the full 99-year term---it turns over meaningfully within the program horizon.

Notably, the savings from the loan structure are \textbf{robust across scenarios} (\$1.15M--\$1.28M), because while conservative assumptions slow repayments, they also increase originations, maintaining a similar proportional benefit.

\subsection{Collateral Adequacy}

The full-principal-recovery assumption deserves explicit justification. At the median property value of \$73{,}896, the \$10{,}000 LEAP cap implies a loan-to-value ratio of just 13.5\%. Even at the 95th percentile of LTV (44.6\%), the property provides more than adequate collateral. Combined with the fact that LEAP liens sit behind existing mortgages and the \$10{,}000 amount is small relative to home equity, the risk of credit loss on any individual loan is negligible. This finding would be worth revisiting only if LEAP were extended to significantly higher loan amounts or to markets with substantially lower property values.

\subsection{Model Calibration}

The model is calibrated to three observable anchors:

\begin{enumerate}[nosep]
  \item \textbf{Eligible population}: 34{,}269 lines from system-materials data.
  \item \textbf{Loan-size distribution}: 64 historical funded amounts, bootstrap-resampled.
  \item \textbf{First-year activity}: 62 active loans observed; model produces 63.7 on average.
\end{enumerate}

The sale-probability assumption (5.6\%) is anchored to the Columbus property sample's average tenure of ${\sim}$18 years. The equity-extraction assumption (2.0\%) is judgmental but consistent with CFPB data on homeowner borrowing behavior~\cite{cfpb2023}.

\subsection{Limitations}

Several limitations should inform the interpretation of results:

\begin{enumerate}[nosep]
  \item \textbf{No street-level replacement schedule.} The city's actual year-by-year replacement plan is not available in the data. We approximate it with a smooth 12-year distribution and test alternatives in sensitivity analysis. Access to the actual schedule would substantially reduce uncertainty.
  \item \textbf{Limited program history.} LEAP has operated for only one observed period (42 completed loans). Long-run uptake behavior---particularly how awareness and contractor capacity evolve---remains uncertain. The model's adoption curve is calibrated to this single observation window and extrapolated forward.
  \item \textbf{No borrower demographics.} Without homeowner age, mortgage status, or direct sale/refinance histories, repayment probabilities rely on aggregate statistical anchors rather than borrower-level models.
  \item \textbf{Neural network targets are heuristic.} The uptake and sale-hazard neural networks are trained on v2-calibrated heuristic labels derived from parcel features, not on observed LEAP enrollment or sale outcomes. They function as learned interpolators of domain assumptions anchored to real v2 program data, providing heterogeneous per-parcel variation while preserving population-level calibration. Their value lies in enabling agent-level heterogeneity, not in independent predictive power. The held-out v4 test set confirms that the learned patterns generalize to unseen synthetic parcel data.
  \item \textbf{Correlation partially addressed.} The base model treats events as independent. The agent-based simulation partially addresses this via housing market dynamics shifting (boom/normal/recession Markov chain), which creates correlated macro shocks across all parcels simultaneously. However, neighborhood-level contagion effects and localized distress clustering are not modeled.
\end{enumerate}

\subsection{Model Validation Summary}

Table~\ref{tab:validation} presents the quantitative validation metrics that assess the internal consistency and precision of the modeling framework.

\begin{table}[H]
  \centering
  \caption{Model validation summary statistics.}
  \label{tab:validation}
  \begin{tabular}{@{}lrl@{}}
    \toprule
    \textbf{Metric} & \textbf{Value} & \textbf{Interpretation} \\
    \midrule
    MC Mean (3{,}000 paths) & \$2{,}489{,}537 & Central estimate \\
    MC P95 (3{,}000 paths) & \$2{,}731{,}561 & Headline recommendation \\
    Deterministic Benchmark & \$2{,}471{,}551 & Closed-form cross-check \\
    Deterministic--MC Gap & $-$0.72\% & Excellent agreement \\
    \midrule
    Bootstrap CI (P95) & [\$2{,}716{,}604, \$2{,}741{,}433] & 95\% confidence interval \\
    Bootstrap CI Width (P95) & \$24{,}828 & Estimation precision \\
    Bootstrap SE (P95) & \$6{,}194 & Standard error of P95 \\
    \midrule
    Coefficient of Variation & 0.11\% & Simulation stability \\
    Simulation Count & 3{,}000 & Paths evaluated \\
    \textbf{Assessment} & \textbf{PASS} & All checks satisfied \\
    \bottomrule
  \end{tabular}
\end{table}

The validation assessment of \textbf{PASS} is based on three criteria: (1)~the deterministic--Monte Carlo gap is below 1\% (actual: 0.72\%), confirming that the simulation engine correctly implements the expected cash-flow logic; (2)~the coefficient of variation is below 1\% (actual: 0.11\%), indicating that 3{,}000 paths provide stable estimates; and (3)~the bootstrap confidence interval for the P95 is narrow relative to the point estimate (${\pm}$0.45\%), confirming sufficient precision for policy decisions.

\subsection{Method Convergence}

A key strength of this analysis is that \textbf{eight independent analytical approaches}---supplemented by stress testing, geographic analysis, value-of-information decomposition, break-even identification, and adaptive funding optimization---converge on consistent conclusions:

\begin{enumerate}[nosep]
  \item \textbf{Cohort-based Monte Carlo simulation} (3{,}000 paths) yields a P95 of \$2.73M with a mean of \$2.49M and a coefficient of variation of just 0.11\%.
  \item \textbf{Deterministic Markov benchmark} produces \$2.47M (within 0.72\% of MC mean), confirming that the simulation engine is correctly computing expected cash flows. The deterministic result falls within expectations given Jensen's inequality for the max-deficit functional.
  \item \textbf{Bootstrap inference} (2{,}000 resamples) shows the P95 estimate is precise to ${\pm}$\$12K (95\% CI width \$24{,}828), meaning 3{,}000 paths are more than sufficient for policy-relevant precision.
  \item \textbf{Tornado sensitivity analysis} (8 parameters, 1{,}000 paths each) confirms the qualitative sensitivity ranking: uptake rate dominates (78.1\% of variance), followed by schedule timing (11.3\%), then repayment speed and inflation (${\sim}$3.7\% each)---consistent with the three-scenario results.
  \item \textbf{NPV analysis} across five discount rates (1\%--5\%) shows the true economic cost is \$2.05M (at 3\%), with total variation of only ${\pm}$3.2\% across the full rate range. The \$2.15M terminal portfolio will continue generating repayments through 2059.
  \item \textbf{Enhanced agent-based simulation} with seven structural extensions (shift in market dynamics, Weibull hazards, urgency dynamics, equity tracking, capacity constraints, fund returns, credit losses) produces a P95 of \$3.60M, establishing an upper-bound envelope that accounts for parcel heterogeneity and macroeconomic correlation.
  \item \textbf{Neural network parcel-level models} (UptakeNet, SaleHazardNet, CostNet), trained on v2 data and evaluated on the v4 test set, achieve validation correlations of 0.997 (uptake) and 0.973 (sale hazard), confirming that the learned patterns generalize to unseen synthetic parcel data.
  \item \textbf{Columbus v4 out-of-sample testing} runs the model with independently-derived inputs from the held-out v4 test set, confirming structural consistency when the input source is substituted entirely.
\end{enumerate}

These eight core methods are further reinforced by five supplementary analyses: (a)~\textbf{stress testing} bounds the tail at \$3.14M--\$5.75M across three named scenarios; (b)~\textbf{geographic hotspot analysis} identifies the spatial concentration of risk; (c)~\textbf{value-of-information analysis} confirms that uptake rate resolution dominates the research priority; (d)~\textbf{break-even analysis} identifies 0.202\% as the critical uptake threshold; and (e)~\textbf{adaptive funding optimization} demonstrates that a staged deployment strategy can match the solvency guarantee while deferring 40\% of capital commitment.

The convergence of these thirteen analytical perspectives---spanning analytical closed-form solutions, stochastic simulations, machine learning, independent data sources, stress scenarios, geographic decomposition, and decision-theoretic optimization---provides strong evidence that the \$2.73M recommendation is robust and not an artifact of any single modeling choice.

%--------------------------------------------------------------
\section{Recommendation}
\label{sec:recommendation}

Based on the analysis, we make the following recommendations:

\textbf{Primary recommendation:} Reserve \textbf{\$2.73 million in 2026} under the loan-based LEAP structure. This amount:

\begin{itemize}[nosep]
  \item Is the \textbf{95th-percentile funding level} from the base Monte Carlo simulation (3{,}000 paths), sufficient in 19 out of 20 simulated futures under base-case assumptions.
  \item Has a \textbf{bootstrap 95\% confidence interval of [\$2.72M, \$2.74M]}, with a CI width of only \$24{,}828---confirming high estimation precision.
  \item Is \textbf{\$1.27 million less} than the equivalent grant-model reserve (\$4.01M), demonstrating the capital efficiency of the deferred-repayment design.
  \item Falls within a \textbf{95\% uncertainty band of \$2.20M to \$2.77M}, indicating that the distribution of outcomes is relatively tight under the base case.
\end{itemize}

\textbf{ML-enhanced upper bound:} The agent-based simulation with neural network predictions, shift in market dynamics, and seven structural enhancements produces a P95 of \textbf{\$3.60M} (mean \$3.17M). This represents a more conservative estimate that accounts for parcel-level heterogeneity, deadline urgency acceleration, and macroeconomic correlation. Decision-makers who wish to incorporate these structural risks should consider reserving in the \textbf{\$2.73M--\$3.60M range}, with the exact level depending on risk tolerance.

\textbf{Conservative planning envelope:} If decision-makers prefer additional cushion, the conservative scenario's 95\%-safe level of \textbf{\$3.49 million} provides a prudent upper bound that accounts for higher uptake, slower repayments, and faster cost inflation.

\textbf{Grant comparison:} Under a grant-only structure (no repayments), the equivalent 95\%-safe reserve rises to \textbf{\$4.01 million} (base case) or \textbf{\$4.77 million} (conservative case). The loan structure's repayment mechanics reduce the City's upfront capital commitment by roughly 31\%.

\subsection{Adaptive Funding Strategy}

Rather than committing the full \$2.73M reserve upfront, the City can adopt a \textbf{two-stage adaptive funding strategy} that achieves the same 95\% solvency guarantee while deferring a portion of the capital commitment. We optimize this strategy by evaluating 560 grid points across initial-reserve and top-up-threshold combinations, selecting the configuration that minimizes expected total capital while maintaining a $\geq$95\% solvency rate.

The optimal strategy is:

\begin{table}[H]
  \centering
  \caption{Optimal adaptive funding strategy (evaluated on 500 paths, 560-point grid).}
  \label{tab:adaptive}
  \begin{tabular}{@{}lr@{}}
    \toprule
    \textbf{Component} & \textbf{Value} \\
    \midrule
    Initial reserve (2026) & \$1{,}638{,}937 \\
    Conditional top-up amount & \$1{,}092{,}625 \\
    Top-up trigger & Cumulative originations exceed threshold by year 3 \\
    Solvency rate achieved & 95.0\% \\
    Top-up probability & 100.0\% \\
    Expected total capital & \$2{,}731{,}561 \\
    \bottomrule
  \end{tabular}
\end{table}

Under this strategy, the City would:
\begin{enumerate}[nosep]
  \item \textbf{Year 1 (2026):} Appropriate \$1.64M as the initial LEAP reserve.
  \item \textbf{By Year 3 (2028):} If cumulative originations exceed the monitoring threshold (see Section~6.2), release the conditional top-up of \$1.09M.
  \item \textbf{Total commitment:} The expected total capital equals \$2.73M---identical to the lump-sum approach.
\end{enumerate}

In the current calibration, the top-up probability is 100\% (all paths trigger the threshold by year 3), so the expected savings versus the lump-sum approach are zero. However, the adaptive framework provides two operational advantages: (1)~it creates a formal \textbf{decision checkpoint} at which the City can reassess uptake trends before committing the second tranche, and (2)~if early program data reveals lower-than-expected uptake, the top-up amount can be reduced based on updated break-even analysis (Section~4.15). The staged approach also aligns with typical municipal budgeting cycles, where multi-year capital commitments are more politically feasible than single-year appropriations.

\subsection{Monitoring Triggers}

To enable real-time fund management, we propose three quantitative \textbf{early-warning indicators} derived from the simulation's distributional output. Each trigger compares an observable KPI against a threshold calibrated to the simulation's percentile distribution.

\begin{table}[H]
  \centering
  \caption{Monitoring triggers for LEAP fund management.}
  \label{tab:monitoring}
  \begin{tabular}{@{}p{3.5cm}ccp{5cm}@{}}
    \toprule
    \textbf{KPI} & \textbf{Threshold} & \textbf{Baseline} & \textbf{Action if Breached} \\
    \midrule
    Cumulative originations by end of 2028 & $>$ 193 (P75) & 183 (median) & Accelerate reserve top-up; uptake is running hotter than modeled. \\
    Average funded loan size & $>$ \$8{,}312 (P75) & \$8{,}252 (median) & Cost inflation is exceeding projections; revise outflow forecasts upward. \\
    Cumulative repayments by end of 2029 & $<$ 23 (P25) & 27 (median) & Repayment assumptions may be optimistic; consider increasing the reserve buffer. \\
    \bottomrule
  \end{tabular}
\end{table}

These triggers are designed to detect \textbf{structural deviations} from the model's assumptions early enough for the City to respond. The origination trigger (193 by 2028) would fire if the LEAP uptake rate is approximately 5\% above the median path, signaling that the program is on a higher-demand trajectory. The loan-size trigger (\$8{,}312) would fire if construction cost inflation is running at 4\%+ rather than the assumed 3\%. The repayment trigger ($<$23 by 2029) would fire if housing market conditions are suppressing sales below the 5.6\% annual rate assumption.

Each trigger is independently actionable: breaching one trigger does not necessarily imply the others will be breached, as they monitor orthogonal risk dimensions (demand volume, cost per unit, and capital recycling speed). The combination of all three triggers provides comprehensive coverage of the model's key assumptions.

%--------------------------------------------------------------
\section{Conclusion}
\label{sec:conclusion}

Columbus's LEAP fund presents a well-defined financial engineering problem: sizing an upfront reserve to bridge uncertain, front-loaded loan originations against delayed, back-loaded repayments. We have addressed this problem with a comprehensive multi-method quantitative framework that substantially exceeds the depth of a single-model analysis.

\textbf{Primary finding} A \textbf{\$2.73 million loan-based reserve} (95th percentile, base case) is sufficient to operate the LEAP fund through 2037, with a bootstrap 95\% confidence interval of [\$2.72M, \$2.74M]. This estimate is validated by eight independent analytical methods and five supplementary analyses, all converging on a consistent range. The deterministic Markov benchmark independently yields \$2.47M (0.72\% below the MC mean), and the simulation achieves a coefficient of variation of just 0.11\%---confirming both accuracy and precision.

\textbf{Capital efficiency} The deferred-repayment loan structure reduces the required reserve by \textbf{\$1.27 million} (31\%) compared to a grant-only alternative (\$4.01M). The loan portfolio has a half-life of 8.7 years, with 42\% of capital recycled by 2037 and projected 80\% cumulative recovery by 2051. In net present value terms at 3\% discount, the true economic cost is \$2.05M---robust to ${\pm}$3.2\% across the 1\%--5\% discount rate range.

\textbf{Risk characterization.} Sensitivity analysis identifies LEAP uptake rate as the dominant risk driver, accounting for \textbf{78.1\%} of funding variance with a \$987K tornado swing. Stress testing bounds the tail across three named scenarios: a 2029 housing crash raises the P95 to \$3.14M (+14.9\%), a surge in uptake to \$4.98M (+82.1\%), and a ``perfect storm'' combination to \$5.75M (+110.7\%). The ML-enhanced agent simulation, incorporating parcel heterogeneity, shift in market dynamics, and six additional structural features, produces an upper-bound P95 of \$3.60M. Geographic analysis identifies zip codes 43201 and 43206 as the highest-risk areas (composite scores 0.92 and 0.84).

\textbf{Decision support.} Beyond the point estimate, this analysis provides actionable tools for fund management: (1)~an adaptive two-stage funding strategy (\$1.64M initial reserve + \$1.09M conditional top-up) that creates a formal reassessment checkpoint; (2)~three quantitative monitoring triggers (cumulative originations by 2028, average loan size, cumulative repayments by 2029) for early detection of structural deviations; (3)~break-even analysis showing 0.202\% as the critical uptake rate; and (4)~value-of-information rankings to guide future data collection priorities.

\textbf{Robustness.} The convergence of thirteen analytical perspectives---spanning closed-form Markov analysis, 3{,}000-path Monte Carlo simulation, 2{,}000-resample bootstrap inference, tornado sensitivity, NPV at five discount rates, enhanced agent-based simulation with seven structural extensions, three neural network models, out-of-sample validation, three stress scenarios, geographic decomposition, value-of-information analysis, break-even identification, and adaptive funding optimization---provides strong evidence that the recommendation is not an artifact of any single modeling choice or assumption. The model's primary sensitivity is to future LEAP uptake---how many homeowners seek replacement before the city's scheduled program reaches them---rather than to repayment mechanics, loan sizes, or discount rate assumptions.

%--------------------------------------------------------------
\appendix

\section{Reproducibility}

\begin{itemize}[nosep]
  \item \textbf{Base simulation engine:} Python 3, standard library only (no external dependencies).
  \item \textbf{Agent-based simulation:} Requires PyTorch, NumPy, and scikit-learn for neural network training and inference.
  \item \textbf{Simulation count:} 3{,}000 paths per scenario (base MC); 1{,}000 paths (agent-based, sensitivity variants).
  \item \textbf{Random seeds:} Base = 20260328, Optimistic = 20260329, Conservative = 20260330.
  \item \textbf{Entry point:} \texttt{python3 -m solution.run\_analysis}
  \item \textbf{Output directory:} \texttt{output/} (reports, charts, CSV, JSON).
  \item \textbf{Trained models:} Saved to \texttt{solution/trained\_models/} (.pt weights, scaler parameters).
  \item \textbf{Charts:} Generated as SVG using hand-coded rendering (no plotting libraries required).
\end{itemize}

\section{Additional Data That Would Improve Accuracy}

\begin{enumerate}[nosep]
  \item A street-by-street scheduled replacement calendar by year.
  \item Loan-level LEAP history with origination date, current status, and any repayment dates.
  \item Borrower age or seniority indicators to better model inheritance timing.
  \item City data on refinance-triggered repayments or lien releases.
  \item Contractor-cost inflation assumptions from recent bid history or city projections.
\end{enumerate}

\section{Markov Chain Derivations}

The absorbing Markov chain formulation in Section~3.5 provides closed-form expressions for all lifecycle properties used in the paper. This appendix derives the key results.

\paragraph{Fundamental matrix.} For the transient sub-matrix $\mathbf{Q}$ with states Active and Inherited:
\[
  \mathbf{I} - \mathbf{Q} =
  \begin{pmatrix} 0.084 & -0.008 \\ 0 & 0.090
  \end{pmatrix}
\]
Inverting the $2 \times 2$ matrix (determinant $= 0.084 \times 0.090 = 0.00756$):
\[
  \mathbf{N} = (\mathbf{I} - \mathbf{Q})^{-1} = \frac{1}{0.00756}
  \begin{pmatrix} 0.090 & 0.008 \\ 0 & 0.084
  \end{pmatrix} =
  \begin{pmatrix} 11.905 & 1.058 \\ 0 & 11.111
  \end{pmatrix}
\]

\paragraph{Expected absorption time.} $\mathbf{t} = \mathbf{N}\mathbf{1}$:
\[
  t_A = 11.905 + 1.058 = 12.963, \quad t_I = 0 + 11.111 = 11.111
\]

\paragraph{Variance.} Using the standard result $\text{Var} = (2\mathbf{N}_{\text{dg}} - \mathbf{I})\mathbf{t} - \mathbf{t}_{\text{sq}}$, where $\mathbf{N}_{\text{dg}}$ is the diagonal matrix of $\mathbf{N}$:
\[
  \text{Var}(T_A) = (2 \times 11.905 - 1) \times 12.963 + (2 \times 1.058) \times 11.111 - 12.963^2 = 151.06
\]
\[
  \sigma_A = \sqrt{151.06} = 12.29 \text{ years}
\]

\paragraph{Cumulative repayment at year $n$.} The state vector after $n$ steps from the initial state $(1, 0, 0)$ is $\boldsymbol{\pi}^{(n)} = (1, 0, 0) \cdot \mathbf{P}^n$. The cumulative repayment probability is $\pi_R^{(n)} = 1 - \pi_A^{(n)} - \pi_I^{(n)}$, which grows monotonically toward 1 as $n \to \infty$ since the Repaid state is absorbing.

\paragraph{Cohort recovery.} For a loan originated in year $Y$, the number of transition steps within the program horizon is $n = 2037 - Y$. The recovery fraction is $\pi_R^{(n)}$, yielding the values in Table~\ref{tab:markov_results}.

%--------------------------------------------------------------
\begin{thebibliography}{99}

  \bibitem{epa2024}
  U.S.\ Environmental Protection Agency, ``Lead and Copper Rule Improvements,''
  \textit{Federal Register}, vol.~89, no.~86, pp.~32532--32700, 2024.

  \bibitem{cornwell2003}
  D.~A.~Cornwell, R.~A.~Brown, and S.~H.~Via, ``National survey of lead service line occurrence,''
  \textit{Journal AWWA}, vol.~108, no.~4, pp.~E182--E191, 2016.

  \bibitem{hanna2012}
  M.~Hanna-Attisha, J.~LaChance, R.~C.~Sadler, and A.~Champney~Schnepp, ``Elevated blood lead levels in children associated with the Flint drinking water crisis: a spatial analysis of risk and public health response,''
  \textit{American Journal of Public Health}, vol.~106, no.~2, pp.~283--290, 2016.

  \bibitem{efron1979}
  B.~Efron, ``Bootstrap methods: another look at the jackknife,''
  \textit{The Annals of Statistics}, vol.~7, no.~1, pp.~1--26, 1979.

  \bibitem{efron1993}
  B.~Efron and R.~J.~Tibshirani, \textit{An Introduction to the Bootstrap}.
  Boca Raton, FL: Chapman \& Hall/CRC, 1993.

  \bibitem{kemeny1960}
  J.~G.~Kemeny and J.~L.~Snell, \textit{Finite Markov Chains}.
  Princeton, NJ: Van Nostrand, 1960.

  \bibitem{grinstead1997}
  C.~M.~Grinstead and J.~L.~Snell, \textit{Introduction to Probability}.
  Providence, RI: American Mathematical Society, 1997.

  \bibitem{rubinstein2016}
  R.~Y.~Rubinstein and D.~P.~Kroese, \textit{Simulation and the Monte Carlo Method}, 3rd~ed.
  Hoboken, NJ: Wiley, 2016.

  \bibitem{glasserman2003}
  P.~Glasserman, \textit{Monte Carlo Methods in Financial Engineering}.
  New York: Springer, 2003.

  \bibitem{lawless2003}
  J.~F.~Lawless, \textit{Statistical Models and Methods for Lifetime Data}, 2nd~ed.
  Hoboken, NJ: Wiley, 2003.

  \bibitem{bonabeau2002}
  E.~Bonabeau, ``Agent-based modeling: methods and techniques for simulating human systems,''
  \textit{Proceedings of the National Academy of Sciences}, vol.~99, suppl.~3, pp.~7280--7287, 2002.

  \bibitem{goodfellow2016}
  I.~Goodfellow, Y.~Bengio, and A.~Courville, \textit{Deep Learning}.
  Cambridge, MA: MIT Press, 2016.

  \bibitem{saltelli2008}
  A.~Saltelli, M.~Ratto, T.~Andres, F.~Campolongo, J.~Cariboni, D.~Gatelli, M.~Saisana, and S.~Tarantola, \textit{Global Sensitivity Analysis: The Primer}.
  Chichester, UK: Wiley, 2008.

  \bibitem{howard1966}
  R.~A.~Howard, ``Information value theory,''
  \textit{IEEE Transactions on Systems Science and Cybernetics}, vol.~2, no.~1, pp.~22--26, 1966.

  \bibitem{cfpb2023}
  Consumer Financial Protection Bureau, ``Data point: mortgage market activity and trends,''
  Office of Research, Washington, D.C., 2023.

  \bibitem{ahs2021}
  U.S.\ Census Bureau, ``American Housing Survey,''
  Washington, D.C., 2021.

  \bibitem{fabozzi2007}
  F.~J.~Fabozzi, \textit{Fixed Income Analysis}, 2nd~ed.
  Hoboken, NJ: Wiley, 2007.

  \bibitem{eschenbach1992}
  T.~G.~Eschenbach, ``Spiderplots versus tornado diagrams for sensitivity analysis,''
  \textit{Interfaces}, vol.~22, no.~6, pp.~40--46, 1992.

  \bibitem{troesken2006}
  W.~Troesken, \textit{The Great Lead Water Pipe Disaster}.
  Cambridge, MA: MIT Press, 2006.

\end{thebibliography}

\end{document}
